\scanpage{571}
\mspage{480}
de la famille (pas néc. petite) de sous-catégories
fermées $ {}_xM$ de $M$ (\unsure{indexée|indexées} par
$x\in\operatorname{Ob}\underline S$), \inserted{accessibles dans $M$ en vertu de 4.16.7.2,} donc notre assertion est un
cas particulier de 4.16.7.1.

On peut aussi, pour toute partie petite
$J$ de $\operatorname{Ob}\underline S$, considérer la sous-catégorie
fermée $\underline S_J$\inserted{$\subset\underline S$} de $\deleted{\unsure{\underline S}}M$ engendrée
par $J$. L'ens \inserted{$I$} des parties $J$ en
question, ordonné par inclusion, est
filtrant devant \deleted{\unsure{tout}} cardinal $\pi\in\unsure{\mathcal U|U}$, et
$\underline S$ est réunion croissante des $\underline S_J$.
D'autre part, la proposition nous dit
que $\underline S$ est accessible dans $M$ ssi
\deleted{\unsure{$J$}} il existe un $J\in I$, tel que $\underline S_J=\underline S$,
i.e. tel que $\underline S_J=\underline S_{J'}$ pour $J\unsure{\subseteq|\subset}J'$.

\uline{Proposition} 4.16.7.4. Soit $M$ une catégorie
accessible. Les conditions suivantes
sont équivalentes

a) Il existe une sous-catégorie $\underline S$ de $M$
fermée dans $M$, et non accessible dans $M$.

b) \deleted{Il existe \unsure{soit|des} \unsure{$\alpha$|objets} \unsure{premier|paramétrés} ordinal} \inserted{Soit $I$ l'ens des ordinaux $\alpha$}
tels que $\operatorname{card}\alpha$ \deleted{\unsure{$\leq$}} $\in\mathcal U$ \deleted{\unsure{; (donc}}\inserted{\unsure{i.e.}} des
% archive sheet number: 570
